\chapter{Conclusiones y Trabajo Futuro}
\label{cap:conclusiones}

\section{Conclusiones}

El desarrollo del Sprint 1 ha permitido materializar con éxito los objetivos planteados al inicio de la memoria, combinando la construcción de una aplicación funcional con la adhesión estricta a un proceso riguroso de ingeniería del software:
\begin{itemize}
    \item \textbf{Objetivo principal:} Se diseñó, implementó y desplegó en producción la arquitectura base de un servicio SaaS con gestión integral de usuarios, automatizando por completo el flujo desde la codificación hasta la entrega en la nube.
    \item \textbf{Diseño arquitectural desacoplado:} Se estructuró el backend en capas bien diferenciadas (API, lógica y datos) servidas desde el mismo origen que el cliente web, garantizando la independencia y modularidad del sistema.
    \item \textbf{Gestión de identidades y sesiones:} Se implementó la autenticación local con derivación de contraseñas mediante \texttt{bcryptjs} y almacenamiento de sesiones persistentes en base de datos resistente a reinicios.
    \item \textbf{Control de acceso basado en roles:} Se implantó la autorización diferenciada para usuarios y administradores, con baja lógica y protección de rutas en el servidor.
    \item \textbf{Persistencia relacional y configuración:} Se integró PostgreSQL en la nube mediante migraciones SQL versionadas y se aislaron los secretos en variables de entorno y Secret Manager.
    \item \textbf{Garantía de calidad mediante pruebas:} Se consolidó una suite de 80 pruebas automatizadas que cubren casos nominales y de error, ejecutables mediante un único comando.
    \item \textbf{Automatización del ciclo de entrega (CI/CD):} Se implantaron flujos de trabajo en GitHub Actions con validación en contenedores Docker y despliegue automático en Google Cloud Run y Render.
    \item \textbf{Procesos ágiles y trazabilidad:} Se gestionó el backlog en DoneTonic cumpliendo la disciplina de GitHub Flow, vinculando cada tarea a su rama y pull request.
\end{itemize}

\section{Trabajo futuro}

De cara a las siguientes fases del proyecto (particularmente el Hito 4), se identifican tres líneas prioritarias de evolución técnica:
\begin{enumerate}
    \item \textbf{Autenticación federada y verificación por correo (P15 y P16):} Incorporar el inicio de sesión con proveedores externos OAuth (Google) y el flujo de confirmación de cuentas mediante un servicio de correo transaccional (activando \texttt{REQUIRE\_EMAIL\_CONFIRMATION=true}).
    \item \textbf{Manejo avanzado de errores y feedback interactivo (P17):} Perfeccionar el sistema de retroalimentación al usuario en la interfaz, unificando los códigos de error y mejorando la recuperación ante fallos de conectividad.
    \item \textbf{Rediseño visual y expansión funcional:} Evolucionar la estética de la interfaz hacia una experiencia comunitaria más rica una vez que se incorporen las funcionalidades de publicación e interacción social.
\end{enumerate}
